\documentclass[11pt]{amsart}
\usepackage[T1]{fontenc}
\usepackage{lmodern,amsmath,amssymb,mathtools,microtype,booktabs}
\usepackage[margin=1.15in]{geometry}
\usepackage[hidelinks]{hyperref}
\newtheorem{theorem}{Theorem}[section]
\newtheorem*{maintheorem}{Main theorem}
\newtheorem*{reciprocaltheorem}{Reciprocal reconstruction theorem}
\newtheorem{lemma}[theorem]{Lemma}
\newtheorem{proposition}[theorem]{Proposition}
\newtheorem{corollary}[theorem]{Corollary}
\theoremstyle{definition}\newtheorem{definition}[theorem]{Definition}
\theoremstyle{remark}\newtheorem{remark}[theorem]{Remark}
\newcommand{\R}{\mathbb R}\newcommand{\Z}{\mathbb Z}
\newcommand{\E}{\mathbb E}\newcommand{\Prob}{\mathbb P}
\newcommand{\dd}{\,\mathrm d}\newcommand{\one}{\mathbf 1}
\DeclareMathOperator{\dist}{dist}\DeclareMathOperator{\covol}{covol}
\DeclareMathOperator{\area}{area}\DeclareMathOperator{\supp}{supp}
\numberwithin{equation}{section}
\title[Scalar collision laws]{Scalar collision laws and recognition of periodic dispersing billiards}
\author{Qian Qi}\date{October 4, 2026}
\subjclass[2020]{37D50, 52A27, 52C23, 62G05}
\keywords{Controlled collision laws, reciprocal preparations, mean exit time,
constructive reconstruction, period recognition, interval enclosures}
\begin{document}\raggedbottom
\begin{abstract}
We reconstruct separated convex obstacles from reciprocal collision
probabilities, with one bit per attempted launch and the same zero
for a solid start and a free miss. The two preparations reveal a
Markov difference of the obstacle occupation. A mean-exit comparison
recovers this occupation without a supplied zero set or a common
positive drift direction. Every nontrivial bounded displacement law
with steps shorter than component separation is admissible; centered
laws give an explicit diameter-to-variance stability bound. A killed
Bellman recursion converges geometrically on a fixed enlarged aperture.
For four pooled directions its updates are finite grid operations,
with exact rational interval enclosures and no direction-resolved
output. On bounded periodic classes with a known positive nonperiod
patch margin and uniform smoothness, threshold clusters, convex hulls
and support smoothing recover the whole table. At localization
$\sigma\asymp\nu^{3/2}$ the construction uses $O(\nu^{-3})$
spatial centers, $O(\nu^{-3}\log(C/(\nu\delta)))$ one-bit attempts
and $O(\nu^{-6})$ arithmetic and fixed-kernel evaluations for
conditional $C^2$ error $O(\nu)$. Calibration and reciprocal input
coupling are explicit assumptions. Exact rational period relations
are distinguished from estimated real coordinates. The deterministic
scalar inverse and all earlier intrinsic-law results are retained;
no uniformly prepared count-law or passive spectral conclusion is
inferred from this active-input experiment.
\end{abstract}\maketitle
\section*{Introduction and main results}
A collision bit is an output, not a description of the entire
experiment. Here the command specifies a spatial probability density,
a direction and a duration. All attempted preparations remain in the
denominator, including an attempted start in the unknown solid. This
convention makes opposed launch laws measure a signed occupation
difference rather than a conditional free-phase probability. The exact
data are functionals of the spatial input; a finite experiment uses a
growing spatial grid. No claim about uniformly prepared count laws is
made by reducing the per-attempt output to a bit.

\begin{maintheorem}
On the bounded periodic smooth class $\mathcal B_\eta$ of
Definitions~\ref{def:class} and \ref{def:margin}, with a known positive
patch margin $\eta$, a finite scalar experiment and a finite geometric
algorithm recover a smooth periodic presentation to $C^2$ error
$O(\nu)$, its primitive orbit count and the exact rational relations
of its period generators, with confidence $1-\delta$.
Each attempt returns one bit: a free segment hit is one; a solid start
and a free miss are both zero, and all attempts remain in the denominator.
The design uses $O(\nu^{-3})$ localized spatial commands of scale
$\sigma\asymp\nu^{3/2}$ and
$O(\nu^{-3}\log(C/(\nu\delta)))$ attempts. In the coupling model of
Theorem~\ref{thm:calibration}, combined position, timing and angular
displacement error at most $c\sigma$ is allowed. The arithmetic bound
is $O(\nu^{-6})$ in the model of Theorem~\ref{thm:constructive}; the
real period basis is estimated, not recovered exactly. All constants
and admissible accuracies depend on the displayed priors and $\eta$.
\end{maintheorem}

The central fixed-displacement identities are
\[
 B_a(q)-B_{-a}(q+a)=\one_{\mathcal O}(q+a)-\one_{\mathcal O}(q),
 \qquad
 u(x)=-\min_{0\le k\le m}\sum_{j<k}\delta_a u(x+ja).
\]
The second identity assumes a zero of the nonnegative $u$ on each chain.
A separation bound and a component-diameter bound provide that zero,
including for smoothed occupation. Sections~\ref{sec:reversal}--\ref{sec:finite}
retain the complete fixed-displacement inverse and its period-recognition
and existence-based finite-observation consequences. The principal
extensions in this article are as follows.

\smallskip\noindent
\emph{Randomized preparation without exact collimation.}
For a known displacement distribution $\rho$, the translated opposed
protocol measures $g=(T_\rho-I)u$. If the associated walk hits
$\{u=0\}$ within $m$ steps, the recursion
\[
 V_0=0,\qquad V_{n+1}=\max(0,T_\rho V_n-g)
\]
recovers $u$ at step $m$, with uniform forcing-error amplification at
most $m$. Theorem~\ref{thm:stopped} supplies a geometric cone condition
ensuring this stopping bound and allows a fixed continuous spread of
nominal directions and durations. This exact extension requires the
specified coupled reverse preparation; it is not a claim that arbitrary
unknown jitter preserves the balance identity.

\smallskip\noindent
\emph{Quantified apparatus tolerance.}
For localized inputs of scale $\sigma$, Theorem~\ref{thm:calibration}
controls the bias under a coupled start-position error $\ell$, timing
error $\tau$ and angle error $\alpha$ by
$C(\ell+\tau+b\alpha)/\sigma$, when the nominal duration is at most
$b$ and the combined error is at most $\sigma$. The estimate includes
grazing directions and errors depending on the commanded start. It is
a geometric stability bound, not a calibration procedure for an apparatus.
In particular, scale-independent propagation of a mean error does not
remove the need to sharpen input calibration with $\sigma$.

\smallskip\noindent
\emph{A finite geometric reconstruction.}
On the bounded periodic class $\mathcal B_\eta$ of
Definitions~\ref{def:class} and \ref{def:margin}, with a \emph{known}
positive nonperiod patch margin $\eta$, Theorem~\ref{thm:constructive}
uses the same fixed-duration opposed scalar commands to construct
complete component approximations by thresholding, a finite proximity
graph and convex hulls. No search over a compact infinite-dimensional
shape class is used in this theorem. Compensated smoothing supplies
$C^2$ supports and the retained finite patch test locks the period
relations. The uniform error, probability and operation bounds depend
on $\eta$ and all stated geometric and smoothness priors. The operation
count is in the declared arithmetic and fixed-kernel evaluation model,
not a bit-complexity or optimality assertion. Exact rational relations
refer to an accepted generator basis whose real coordinates have small
error, not to exact recovery of arbitrary real numbers.

These conclusions leave the earlier existence regularization in
Theorem~\ref{thm:finite} unchanged and identify a separate constructive
route. They do not remove the bounded periodic-presentation hypothesis
or certify an unknown patch margin at a finite time. The exact local
inverse itself does not require periodicity. The closest observation
framework is hit probability and active geometric probing; the precise
capacity-functional relation and theorem-level comparisons are in
Section~\ref{sec:hit-comparison}.

\subsection*{Reciprocal randomization and the main extension}
The deterministic and cone-directed results above remain valid. The
following extension removes the almost-sure directional progress
condition, not the controlled spatial input or all-attempt convention.

\begin{reciprocaltheorem}
For separated components of bounded diameter, two reciprocal spatial-input
collision functionals determine the obstacle union whenever the known
random displacement is nontrivial, bounded, and shorter than the
component separation. They output pooled hit probabilities, not a
collision position or a direction-resolved mean. Solid starts and free
misses remain zero in the common attempted-preparation denominator.
For centered increments with second moment at least $v_*>0$, the
mollified occupation inverse obeys
\[
 \|u-\widetilde u\|_\infty\le
       \frac{(D+b)^2}{v_*}\,
       \|(T_\rho-I)(u-\widetilde u)\|_\infty,
\]
where $D$ bounds the enlarged component diameters and $b$ the step
lengths. The unknown zero sets need not agree. A Bellman iteration
converges with a uniform geometric tail, including on a fixed observation
aperture enlarged by a prior-controlled collar. No free reference point
or stopping-set measurement is supplied.

For the four-direction pooled protocol of
Theorem~\ref{thm:pooled-recovery}, on the bounded periodic class
$\mathcal B_\eta$ with known positive patch margin $\eta$, a finite
fixed-aperture algorithm has conditional $C^2$ error $O(\nu)$ with
confidence $1-\delta$. It uses $O(\nu^{-3})$ spatial centers,
$O(\nu^{-3}\log(C/(\nu\delta)))$ attempts and $O(\nu^{-6})$
arithmetic and fixed-kernel evaluations at localization
$\sigma\asymp\nu^{3/2}$. Exact rational interval Bellman updates
certify the local occupation bounds on the simultaneous mean-accuracy
event. Calibration, smoothness, period-presentation bounds and the
known $\eta$ are substantive priors. Euclidean basis coordinates are
estimated; only their bounded-denominator relations are identified
exactly. No minimax or Turing bit-complexity statement is asserted.
\end{reciprocaltheorem}

The exact result is proved in Section~\ref{sec:mean-exit}. Its constants
come from exit-time estimates, rather than convergence of an increasing
jet recursion or a global elliptic inverse with known boundary values.
Section~\ref{sec:pooled-finite} implements a finite-support law by exact
grid shifts and rational intervals. General continuous laws are included
in the exact functional theorem, not silently treated as finitely
integrable data. The classical martingale, Green and optimal-stopping
arguments are credited separately from their use with the measured
collision forcing and its unknown zero set.


\section{Reversal balance and a finite-range inverse}
\label{sec:reversal}
A spatial first-impact law may reveal the boundary directly through its
support. We ask instead what remains when the output of each attempt is
only a collision bit. The starting distribution is now a controlled
input. The distinction matters: reducing the recorded output while
allowing localized preparations does not establish dominance over a
uniform-launch experiment. Our principal identity separates the two
issues. It expresses an occupation difference by two scalar collision
probabilities at the same positive horizon, with no small-time limit.

Let $\mathcal O\subset\R^d$ be closed, and put $\chi=\one_{\mathcal O}$.
An attempt starts at $q$ with velocity $v$, $|v|=1$, and runs until time
$t>0$. A start in $\mathcal O$ is an unsuccessful preparation and is
assigned output zero, as is a free start with no collision by time $t$.
A free start with a collision by that time has output one. The two zero
outcomes are not distinguished by the observation. The event includes
the endpoint time; boundary conventions are immaterial for the smooth
classes and absolutely continuous preparations below. Write $a=tv$ and
\begin{equation}\label{eq:bit}
 B_a(q)=\one_{\{q\notin\mathcal O,\ [q,q+a]\cap\mathcal O\ne\varnothing\}}.
\end{equation}
The straight segment in this definition is the unreflected path. It
detects the physical first collision, because the trajectory is straight
until that collision; no part of the path after the first collision is
observed. For a nonnegative probability density $f$ set
\[
 P_a(f)=\int_{\R^d}f(q)B_a(q)\dd q,\qquad f_a(q)=f(q-a).
\]
The denominator is the number of attempted preparations, not the number
of free or successful starts. The experiment commanding $f_a$ shifts
the starting distribution by $+a$ and reverses the velocity. It does
not run a time-reversed reflected trajectory.

\begin{theorem}[Reversal balance]\label{thm:balance}
For every closed obstacle set and every displacement $a$,
\begin{align}
 B_a(q)-B_{-a}(q+a)&=\chi(q+a)-\chi(q),\label{eq:balance-point}\\
 P_a(f)-P_{-a}(f_a)&=
       \int f(q)\{\chi(q+a)-\chi(q)\}\dd q.\label{eq:balance}
\end{align}
The pointwise identity requires no convexity, separation or periodicity.
It remains valid when the segment intersects several components.
\end{theorem}
\begin{proof}
Let $I(q)$ indicate whether $[q,q+a]$ meets $\mathcal O$. The two
oriented segments have the same image, so the left side of
\eqref{eq:balance-point} is
$\{\chi(q+a)-\chi(q)\}I(q)$. If either endpoint is in $\mathcal O$,
then $I(q)=1$; if neither is, the prefactor is zero. This proves the
pointwise assertion. In the reverse probability change variables from
$q+a$ to $q$. The result is \eqref{eq:balance}.
\end{proof}

In particular, the pair of functionals $P_a,P_{-a}$ on smooth compactly
supported launch densities determines the signed measure with density
$\delta_a\chi(q):=\chi(q+a)-\chi(q)$. Testing all nonnegative smooth
densities determines this measure uniquely, by scaling and linearity.
These are functionals of a controllable spatial input, not two real
observations. For finite acquisition we use an explicit finite family
of inputs rather than appeal to exact differentiation of such a measure.

A difference operator ordinarily leaves an additive constant along each
translation orbit undetermined. The next theorem fixes it by a geometric
zero condition, without providing a free reference point.

\begin{theorem}[Finite-chain inversion]\label{thm:chain}
Let $u\ge0$ be a function, $a\ne0$, and $m\ge1$. Suppose that for every
$x$ in a region $U$,
\begin{equation}\label{eq:zero-chain}
                 \min_{0\le k\le m}u(x+ka)=0.
\end{equation}
Then the values of $\delta_a u$ on
$U^+=\bigcup_{j=0}^{m-1}(U+ja)$ determine $u$ on $U$ by
\begin{equation}\label{eq:min-inverse}
 u(x)=-\min_{0\le k\le m}\sum_{j=0}^{k-1}\delta_a u(x+ja),
\end{equation}
where the empty sum is zero. If $u'$ obeys the same zero condition,
then, for $1\le p\le\infty$,
\begin{equation}\label{eq:chain-stability}
 \|u-u'\|_{L^p(U)}\le
       m\|\delta_a u-\delta_a u'\|_{L^p(U^+)}.
\end{equation}
More generally, perturbing every difference in a chain by at most $b$
perturbs the reconstruction by at most $mb$.
\end{theorem}
\begin{proof}
The partial sum is $u(x+ka)-u(x)$. Its minimum equals $-u(x)$ by
\eqref{eq:zero-chain}, proving the formula. The minimum of two finite
lists differs by at most their maximum entrywise difference. Hence
\[
 |u(x)-u'(x)|\le
   \sum_{j=0}^{m-1}|\delta_a u(x+ja)-\delta_a u'(x+ja)|.
\]
Minkowski's inequality and translation of each integration domain prove
\eqref{eq:chain-stability}; the uniform perturbation statement follows
from the same inequality. Almost-everywhere hypotheses suffice: only
finitely many translates of exceptional null sets occur.
\end{proof}

\begin{lemma}[A geometric zero witness]\label{lem:zero}
Suppose $\mathcal O$ is a union of compact components of diameter at
most $D_0$, with mutual distance at least $d_0>0$. If
$0<t<d_0$, $a=tv$ and $mt>D_0$, then $u=\chi$ satisfies
\eqref{eq:zero-chain} everywhere.
Fix $\sigma_0>0$ with $t<d_0-2\sigma_0$ and
$mt>D_0+2\sigma_0$. If $\kappa\ge0$ has integral one and support
in the closed unit ball, then for every $0<\sigma\le\sigma_0$ the
local occupation average
\begin{equation}\label{eq:occupation}
 u_\sigma(x)=\int \sigma^{-d}\kappa((q-x)/\sigma)\chi(q)\dd q
\end{equation}
also satisfies \eqref{eq:zero-chain}.
\end{lemma}
\begin{proof}
If all chain points belong to $\mathcal O$, successive points cannot
belong to different components, since their distance is $t<d_0$.
They therefore lie in one component, contradicting the endpoint distance
$mt>D_0$. For the second assertion, $u_\sigma(x)>0$ implies
$x\in\mathcal O+\sigma\overline B$. These enlarged components have
diameter at most $D_0+2\sigma$ and separation at least $d_0-2\sigma$.
The same argument excludes positivity at every point of the chain.
\end{proof}

\begin{corollary}[Local rigidity from scalar launch laws]\label{cor:local}
Under the diameter and separation hypotheses, the two probability
functionals at one fixed horizon and opposite directions determine the
obstacle indicator almost everywhere on $U$, using only launch densities
supported on its bounded chain enlargement. For regular closed bodies
they determine the bodies themselves in the interior of $U$.
Alternatively, the same determination follows from localized kernel
preparations at all sufficiently small scales, with the explicit
finite-chain formula at each scale.
\end{corollary}
\begin{proof}
Apply Theorem~\ref{thm:balance} to recover $\delta_a\chi$ and
Theorem~\ref{thm:chain} with Lemma~\ref{lem:zero}. Two different regular
closed convex unions would differ on an open set, so almost-everywhere
equality determines the set. For the alternative, the balance identity
recovers $\delta_a u_\sigma$ exactly, the minimum formula recovers
$u_\sigma$, and approximate-identity convergence gives $\chi$ almost
everywhere as $\sigma\downarrow0$.
\end{proof}

The zero witness is substantive. For constant sequences both zero and
one have zero difference, and for unbounded slabs parallel to $a$ the
same ambiguity can persist. Here the witness is implied by bounds on
physical component size and separation, not by an unobserved boundary
value. The local theorem uses no periodicity or boundary smoothness.
In the convex billiard class and at $t<d_0$, there cannot be a second
collision: a reflected exterior ray cannot re-enter the same convex
body, and another body is at least $d_0$ away. Thus the successful
output is also the one-collision count at this horizon, with solid
preparations still counted as zero attempts.

\section{Finite scalar acquisition of a complete patch}
\label{sec:acquisition}
We now work in the plane. Fix once and for all a nonnegative, even
$C^\infty$ probability kernel $\kappa$ supported in the unit disk, and
put $K=\|\nabla\kappa\|_1$. We may increase $K$ to be at least one.
Take, for example,
\begin{equation}\label{eq:design}
 t=\min(1,d_0/4),\qquad \sigma_0=\min(1,d_0/8),\qquad
 m=\left\lfloor\frac{D_0+2\sigma_0}{t}\right\rfloor+1,
 \qquad a=t(1,0).
\end{equation}
The chain is fixed by the numerical priors, not by an unknown boundary.
For $f_{x,\sigma}(q)=\sigma^{-2}\kappa((q-x)/\sigma)$ define
\[
 d_\sigma(x)=P_a(f_{x,\sigma})-P_{-a}(f_{x+a,\sigma}).
\]
Equation~\eqref{eq:balance} and Lemma~\ref{lem:zero} give
\begin{equation}\label{eq:smoothed-min}
 d_\sigma(x)=u_\sigma(x+a)-u_\sigma(x),\qquad
 u_\sigma(x)=-\min_{0\le k\le m}\sum_{j<k}d_\sigma(x+ja).
\end{equation}
There is no deconvolution, time derivative or noise amplification by
$\sigma^{-1}$ in this step.

\begin{proposition}[Stable recovery of occupation averages]
\label{prop:mean-stability}
At each point $x$ the $2m$ scalar means in \eqref{eq:smoothed-min}
determine $u_\sigma(x)$. If each mean has absolute error at most
$\epsilon$, the error is at most $2m\epsilon$, uniformly in $\sigma$.
Clipping the answer to $[0,1]$ cannot increase this error. For two
physical candidates whose corresponding means differ by at most $b$,
their occupation averages at $x$ differ by at most $2mb$.
\end{proposition}
\begin{proof}
Each balance difference has error at most $2\epsilon$, or at most
$2b$ in the comparison version. Apply the perturbation conclusion of
Theorem~\ref{thm:chain}. The true average lies in $[0,1]$.
\end{proof}

Each mean is estimated by repeated attempts with the prescribed input
law. The actual start position need not be recorded with the output;
the commanded distribution, its center and the direction are known.
The only per-attempt recorded variable is the bit. If the actual launch
law differs from its command by total variation at most $\zeta$, the
bias of every mean is at most $\zeta$. In particular a center error
$\ell$ changes this smooth input density in $L^1$ by at most
$K\ell/\sigma$. These observations describe input calibration costs;
they are not a claim that arbitrary directional or dynamical apparatus
errors are harmless.

\begin{lemma}[A finite grid controls the patch geometry]
\label{lem:grid}
Assume all physical bodies are strictly convex, have curvature at most
$\kappa_+$, and satisfy the common separation and diameter bounds.
Let $U$ be a square containing, with a collar of width $2D_0+2$, every
body to be compared. Take
$\sigma<\min(\sigma_0,(2\kappa_+)^{-1},d_0/20,1/4)$.
A grid of covering radius $\sigma/(32K)$ in $U$ contains
$O_U(\sigma^{-2})$ points. If the scalar means for two admissible
physical candidates differ at every grid command by at most
$b=1/(32m)$, then their complete bodies in the protected inner patch
can be matched with Hausdorff error at most $2\sigma$.
\end{lemma}
\begin{proof}
By Proposition~\ref{prop:mean-stability} the occupation averages differ
by at most $1/16$ at each grid point. Each occupation function has
Lipschitz constant at most $K/\sigma$. The chosen covering radius
therefore gives a uniform difference at most $1/8$ on $U$.

For a convex body $C$, let $C_{-\sigma}$ denote its erosion by the
closed radius-$\sigma$ disk. If $x\in C_{-\sigma}$, then
$u_\sigma(x)=1$, whereas if $\dist(x,\mathcal O')>\sigma$, the other
occupation average is zero. It follows that
\begin{equation}\label{eq:erosion-inclusion}
 \mathcal O_{-\sigma}\cap U\subset\mathcal O'+\sigma\overline B,
 \qquad
 \mathcal O'_{-\sigma}\cap U\subset\mathcal O+\sigma\overline B.
\end{equation}
The notation on the left means the union of the component erosions.
The curvature upper bound ensures an interior rolling radius
$r=1/\kappa_+$. One elementary justification uses the support function:
$p+p''\ge r$, so $p-r$ is a convex support function and
$C=C_0+r\overline B$. Hence
$C=(C_{-\sigma})+\sigma\overline B$ when $\sigma<r$.
Each erosion is nonempty and connected.

Distinct expanded target components have positive separation. A
connected erosion in \eqref{eq:erosion-inclusion} must therefore lie
in the radius-$\sigma$ enlargement of a single target body. Restoring
the disk shows $C\subset D+2\sigma\overline B$. Apply the other
inclusion to $D$ to obtain $D\subset E+2\sigma\overline B$ for a
single source body $E$. All these bodies and their erosions are in $U$
by the protective collar. Since $C\subset E+4\sigma\overline B$ and
$4\sigma<d_0$, necessarily $E=C$. Thus $d_H(C,D)\le2\sigma$.
The same argument in the reverse direction makes the matching bijective
for all complete bodies used in the protected patch. No incomplete
outer fragment is declared to be a physical component.
\end{proof}

For later use we record the smooth interpolation underlying the
$C^2$ conclusion. Here $C^j$ norms of support functions are on the
normal-angle circle, in a fixed laboratory frame.
\begin{lemma}[From Hausdorff error to smooth support error]
\label{lem:smoothing}
For convex bodies with bounded $C^6$ support norms, Hausdorff distance
at most $e\le1$ implies $C^2$ support difference at most $C e^{2/3}$.
The same estimate holds for centered support functions, with a changed
constant and their Steiner-center error bounded by $2e$.
\end{lemma}
\begin{proof}
Let $\varphi$ be a smooth even probability kernel of compact support
on the line, and set $\varphi_2(s)=\varphi(s/2)/2$ and
$J=(4\varphi-\varphi_2)/3$. Then $J$ has integral one and its moments
of degrees one, two and three vanish. Periodized convolution at scale
$h$ satisfies
\[
 \|J_h*(p-p')\|_{C^2}\le Ceh^{-2},\qquad
 \|J_h*p-p\|_{C^2}+\|J_h*p'-p'\|_{C^2}\le C h^4.
\]
The first inequality uses $\|p-p'\|_\infty=d_H(C,D)$ and the $L^1$
norms of the first two derivatives of the kernel. The second is Taylor's
formula for $p,p',p''$ through the cancelled orders, using the $C^6$
bound. Choose $h=e^{1/6}$. The Steiner point is
$\pi^{-1}\int_0^{2\pi}p(\theta)(\cos\theta,\sin\theta)\dd\theta$;
its change has norm at most $2e$. Subtracting this translation gives
the centered estimate.
\end{proof}

An admissible-candidate formulation will suffice below. It avoids
claiming an efficient reconstruction algorithm for all smooth shapes:
we select a physical candidate consistent with the finitely many scalar
means, then apply the explicit finite period tests to its recovered
patch. The statements concern both this selection and its quantified
error, not the execution time of searching an infinite-dimensional class.

\section{Recognition of the intrinsic period lattice}
\label{sec:periods}
The scalar inverse recovers a physical patch before periodicity is used.
We next give the local period argument needed to extend that patch.
This argument is retained from the preceding repeated-motif revision
\cite{Supplement}; it is included here so that the primary theorem does
not depend on a different sensor's acquisition proof.

\begin{definition}[Bounded periodic class]\label{def:class}
Fix $r_0,L_0,V_0,D_0,d_0,\kappa_-,\kappa_+>0$. The class
$\mathcal B$ consists of nonempty planar unions
\begin{equation}\label{eq:period-presentation}
 \mathcal O=\bigcup_{i=1}^r(C_i+\Lambda_0),\qquad
 \Lambda_0=L\Z^2,\qquad 1\le r\le r_0,
\end{equation}
with exactly $r$ component orbits modulo $\Lambda_0$,
$\|L\|,\|L^{-1}\|\le L_0$, and $\covol\Lambda_0\ge V_0$.
Every physical component is a compact strictly convex body of diameter
at most $D_0$, with curvature in $[\kappa_-,\kappa_+]$; distinct
components are separated by at least $d_0$. Centered support functions
have $C^{6,\beta}$ norm at most $M_6$, for a fixed $0<\beta\le1$.
The unknown presentation need not be primitive. Repeated shapes and
arbitrary symmetries are allowed. Its full period group is
$\Pi(\mathcal O)=\{w:\mathcal O+w=\mathcal O\}$.
\end{definition}

The high-order bound will only enter finite $C^2$ reconstruction.
Set $B=1+L_0$ and $Q_s=[-s,s]^2$. A convenient fixed observation and
launch aperture is
\begin{equation}\label{eq:aperture}
 R_0=12B+3D_0+4,\qquad R=R_0+mt+2,\qquad Q_R,
\end{equation}
where $t,m$ are from \eqref{eq:design}. Scalar preparations centered
on $Q_{R_0}$ and on the forward chains needed for it are supported in
$Q_R$ for every $\sigma\le\sigma_0$. Each physical flight has the same
fixed duration $t$. There is no recentering at an estimated obstacle.

For a component $C$ write $z_C$ for its Steiner point and $p_C$ for its
centered support function. In the laboratory directions define
\begin{align}
 d_*(C,D)&=\|z_C-z_D\|_\infty+\|p_C-p_D\|_\infty,\label{eq:bodymetric}\\
 \mathcal D_{\mathcal O}(w)&=
 \max_{C:z_C\in Q_B}\ \max_{s\in\{-1,1\}}
                   \inf_{D\in\mathcal C}d_*(C+sw,D),\label{eq:defect}
\end{align}
where $\mathcal C$ is the set of physical components. The metric does
not optimize over rotations. Hausdorff error $e$ implies center error
at most $2e$ and centered-support error at most $3e$. All relevant
infima are attained, since distant centers cannot minimize them.

\begin{lemma}[A finite test for a global period]\label{lem:period}
One has $\mathcal D_{\mathcal O}(w)=0$ if and only if
$w\in\Pi(\mathcal O)$. For $\|w\|_\infty\le8B$ the zero test uses
only components with centers in $Q_B$ and $Q_{9B}$.
\end{lemma}
\begin{proof}
Every $\Lambda_0$-component orbit has a representative with center in
$L[-1/2,1/2]^2\subset Q_B$. If the defect vanishes, each such
representative has both translates by $+w$ and $-w$ among the physical
components. Translating by $\Lambda_0$ gives
$\mathcal O+w\subset\mathcal O$ and $\mathcal O-w\subset\mathcal O$.
The latter implies the reverse of the former inclusion. The converse
is immediate. For the bounded zero test, every possible equal target
has center $z_C\pm w\in Q_{9B}$.
\end{proof}

\begin{lemma}[A protected generating list]\label{lem:generators}
The group $\Pi=\Pi(\mathcal O)$ is a lattice containing $\Lambda_0$,
\begin{equation}\label{eq:index}
 [\Pi:\Lambda_0]\le r\le r_0,\qquad \covol\Pi\ge V_*:=V_0/r_0.
\end{equation}
For any root component centered in $Q_{2B}$, take its center differences
to components centered in $Q_{4B}$ and retain those passing
Lemma~\ref{lem:period}. Their integer span is exactly $\Pi$.
\end{lemma}
\begin{proof}
Map a period coset modulo $\Lambda_0$ to the $\Lambda_0$-orbit of the
translated root. This is injective, since a compact body has no nonzero
translation symmetry. There are at most $r$ such cosets. A finite
extension of a full-rank lattice is a full-rank lattice; the covolume
formula proves \eqref{eq:index}. This quotient group acts freely on
the $r$ component orbits, so its order divides $r$.

All accepted differences are true periods. Conversely, the two columns
of $L$ and representatives of all cosets of $\Pi/\Lambda_0$ can be
chosen to have norm less than $B$, using the same centered fundamental
parallelogram. Translating the root by any of these vectors gives a
target in $Q_{3B}$, strictly inside the cutoff $Q_{4B}$. The accepted
list contains generators of $\Lambda_0$ and representatives of every
coset, and hence generates $\Pi$.
\end{proof}

\begin{theorem}[Whole-table rigidity from scalar collision laws]
\label{thm:global-exact}
On $\mathcal B$, the collision-probability functionals for the two
opposed directions at the fixed duration \eqref{eq:design}, with
controlled smooth starting densities in $Q_R$, determine the entire
obstacle union, its full period lattice $\Pi$, its primitive number
$r_*$ of component orbits, and primitive free area $A_*$. The same is
true for the family of localized kernel commands at all sufficiently
small scales. No impact position, impact angle, collision time, shape
label, integer copy label, primitive cell or free-area normalization
is an observation. The laboratory frame is supplied by the launch
controls; forgetting it leaves a common Euclidean-isometry ambiguity.
\end{theorem}
\begin{proof}
Corollary~\ref{cor:local} recovers the occupation indicator on
$Q_{R_0}$, and regular closedness identifies the physical set there.
Every component with center in $Q_{10B}$, and every component needed
to recognize one of these, lies wholly within the protective collar
of this square. Physical separation identifies the complete components.
Lemmas~\ref{lem:period} and \ref{lem:generators} then recover $\Pi$.

All $\Pi$-component orbits have representatives in $Q_B$. Among
components centered in $Q_{2B}$, two are in the same orbit exactly when
their center difference passes the period test. Indeed a period maps
the first to a physical component with the second's Steiner point;
two disjoint full-dimensional convex bodies cannot have the same
Steiner point, which lies in their interiors. Thus the relation gives
$r_*$ and one representative per orbit. These and $\Pi$ determine the
union everywhere. Finally
\begin{equation}\label{eq:area}
 A_*=\covol\Pi-\sum_{j=1}^{r_*}\area(C_j).
\end{equation}
The conclusion concerns the intrinsic union, not a chosen multiple-cell
presentation. Basis computation from finite exact period vectors is
justified quantitatively in Lemma~\ref{lem:locking} below.
\end{proof}

\begin{definition}[Nonperiod patch margin]\label{def:margin}
For $\eta>0$, let $\mathcal B_\eta$ consist of those members of
$\mathcal B$ such that
\begin{equation}\label{eq:margin}
 w=z_D-z_C,\quad z_C\in Q_{3B},\ z_D\in Q_{5B},\quad
 w\notin\Pi(\mathcal O)\quad\Longrightarrow\quad
                  \mathcal D_{\mathcal O}(w)\ge\eta.
\end{equation}
This concerns translations of a whole central patch, not separation
between individual body shapes. For each fixed table it holds with some
positive $\eta$, because there are finitely many candidate pairs and
Lemma~\ref{lem:period} makes every nonperiod defect positive.
\end{definition}

The next deterministic statement is the link between geometric estimation
and exact discrete information. It does not suppose that perturbed real
vectors themselves generate a discrete group.
\begin{lemma}[Robust period classification and rational locking]
\label{lem:locking}
Fix the numerical bounds of $\mathcal B_\eta$. Suppose complete
components on the protected patch are given by matched approximations
of Hausdorff error at most $e$. For sufficiently small $e$, depending
on $\eta$ and the geometric bounds, finite comparisons recover the true
primitive orbit count and the exact rational coordinates of an accepted
period-generating list in an independent-pair basis. They give a matched
basis for $\Pi$ with error $O(e)$ and an area estimate with error $O(e)$.
\end{lemma}
\begin{proof}
Use approximate centers in $Q_{2B}$ as sources, choose a root there,
and use its targets in $Q_{4B}$. A measured difference $\widehat w$
has a corresponding true center difference $w$ with error at most
$4e$. Test both signs of this translation on all selected sources,
minimizing $d_*$ over approximate components centered in $Q_{10B}$.
For each source-target comparison, the center term changes by at most
$8e$ and the centered-support term by at most $6e$.

If $w$ is a period, all needed partners lie within the target cutoff
by the strict aperture margins, and the measured defect is at most
$14e$. If it is not, its true endpoints lie in $Q_{3B},Q_{5B}$, and
\eqref{eq:margin} applies. Every witness with true center in $Q_B$ is
selected. Restricting the targets cannot reduce a true positive infimum,
so the measured defect is at least $\eta-14e$. Reserve another $6e$
for numerical enclosures of centers and support suprema. Threshold
$\eta/2$, with $20e<\eta/4$, accepts exactly the true periods among
these candidates. The same comparison on all pairs of selected sources
gives their true orbit equivalence relation and hence $r_*$. Targets
for these tests also lie within $Q_{10B}$. Lemma~\ref{lem:generators}
and its strict cutoff margin ensure all protected generators are present.

All accepted vectors have norm at most $U=12B$ for small $e$. A nonzero
determinant of two true periods is an integer multiple of
$\covol\Pi\ge V_*$. Determinant perturbations are $O(Ue)$; after
shrinking $e$, independence is decided at threshold $V_*/2$. Choose
an independent accepted pair with true matrix $D$. The sublattice
$D\Z^2$ has index
\[
       n=|\det D|/\covol\Pi\le Q:=\max(1,\lceil U^2/V_*\rceil).
\]
Every coordinate of $D^{-1}w_j$ has reduced denominator dividing $n$,
and a prior-bounded numerator. The inverse-matrix identity bounds its
estimated error by $C_*e$. Distinct rationals of denominator at most
$Q$ differ by at least $Q^{-2}$. If $C_*e<1/(3Q^2)$, finite search
therefore identifies each exact coordinate uniquely.

Let $q_*$ be their common denominator; it divides $n$ and is at most
$Q$. A column Hermite basis $H$ of the integer span of $q_*I_2$ and
$q_*D^{-1}w_j$ gives
\begin{equation}\label{eq:hermite}
                  \Pi=(DH/q_*)\Z^2.
\end{equation}
The integer group contains $q_*\Z^2$, so the Hermite entries are
bounded in terms of $Q$. Replacing $D$ by its measured version gives
a matched basis with error $O(e)$. This basis need not be a continuously
chosen reduced basis. The parallel-body bound gives an $O(e)$ area
error for each matched convex component; use \eqref{eq:area} for
$A_*$. Certified finite angular meshes and the uniform support Lipschitz
bound enclose all suprema and center integrals to the reserved precision.
\end{proof}

\section{Finite confidence from localized scalar preparations}
\label{sec:finite}
The constants below may depend on the fixed kernel and every numerical
bound of $\mathcal B_\eta$, including $\eta$. We count independent
attempted launches. Spatial preparation precision, deterministic search,
arithmetic and apparatus travel are different resources.

\begin{theorem}[A finite scalar experiment on a patch-margin class]
\label{thm:finite}
On $\mathcal B_\eta$ there are $C,c,\nu_0>0$ such that, for
$0<\nu<\nu_0$ and $0<\delta<1/4$, a deterministic finite list of
localized launch distributions in the fixed square $Q_R$, with only
the two directions $\pm(1,0)$ and the fixed duration $t$, yields the
following guarantee. With probability at least $1-\delta$, it recovers
the primitive orbit count and exact rational relations of a generating
list for $\Pi$, and gives a finite presentation with matched $C^2$
support error, matched period-basis error and primitive-area error at
most $C\nu$. The list uses localization scale
\begin{equation}\label{eq:scalar-scale}
                         \sigma=c\nu^{3/2}
\end{equation}
and at most
\begin{equation}\label{eq:scalar-cost}
                    C\nu^{-3}\log\frac{C}{\nu\delta}
\end{equation}
independent attempts, including solid starts and misses. Only one bit
is recorded per attempt. No output collision localization is required.
A sufficiently small fixed bound on each command's probability bias
is allowed, as specified in the proof.

The uniform exact discrete decisions require the known patch margin
$\eta$. The rate is a conditional sufficient upper bound, not a minimax
rate or an end-to-end computation bound. Localized spatial preparation
and collimated launch directions are substantive inputs.
\end{theorem}
\begin{proof}
\emph{A finite deterministic experiment.}
Take the grid of Lemma~\ref{lem:grid} on $Q_{R_0}$, together with
all its commanded chain centers. Its number of scalar means is
$J\le C m\sigma^{-2}$. At each command estimate its Bernoulli mean
to statistical error at most $\epsilon_0/2$, where
$\epsilon_0=1/(256m)$ is fixed. Allow a systematic mean bias at most
$\epsilon_0/2$. Then each reported mean is within $\epsilon_0$ of
the ideal mean. Hoeffding's inequality and a union bound make this
simultaneous event have probability at least $1-\delta$ with
$C\epsilon_0^{-2}\log(2J/\delta)$ independent attempts per command.
The experiment is deterministic before sampling. It uses no separate
pilot or estimated onset. The total is
\begin{equation}\label{eq:raw-cost}
                      C\sigma^{-2}\log(C/(\sigma\delta)).
\end{equation}

\emph{Selection of an admissible patch.}
The bounded periodic class is precompact in local $C^2$ geometry, and
may be replaced by its closure in a weaker smooth topology retaining
the displayed bounds. To see this, choose the bounded presentations,
pass to a subsequence with fixed $r$, choose component centers in a
closed bounded fundamental parallelogram, and use compactness of the
matrices and support functions. Positive separation and curvature
persist. Higher derivative bounds persist in their usual weak or
H\"older-compact sense. This description also gives a separable
parameter space; a closed compact enlargement satisfying the same
bounds suffices throughout. We do not require the patch-margin subclass
to be closed, since the candidates may be selected in $\mathcal B$.

For each of the finitely many commands the ideal probability is
continuous on this physical class. Except on a null set of initial
points, whether a free segment first meets a strictly convex component
before its endpoint is unchanged under sufficiently small geometric
perturbations. The excluded points lie on body boundaries, translated
endpoint boundaries or tangent lines. Only finitely many bodies can
meet the bounded launch-and-flight region. Dominated convergence proves
continuity of the integral against the fixed smooth command density.

Minimize the maximum mean discrepancy over this class. One can choose
a Borel approximate minimizer within $\epsilon_0$ of the infimum from
a fixed countable dense family: order the family and take its first
eligible member. The infimum of countably many continuous functions
of the data is measurable, so the rule is measurable. On the simultaneous
mean event the true table has discrepancy at most $\epsilon_0$;
the chosen candidate has discrepancy at most $2\epsilon_0$. Its ideal
means differ from those of the truth by at most $3\epsilon_0<1/(32m)$.
Lemma~\ref{lem:grid} gives matched Hausdorff error $e=2\sigma$ for
all complete components used in the central period tests. This step
is an existence-based regularization, not a polynomial-time search
algorithm for arbitrary smooth obstacles.

\emph{Discrete and smooth recovery.}
Use those candidate components as the geometric approximations in
Lemma~\ref{lem:locking}, applying the known margin of the true table.
The candidate itself need not have the same exact symmetry group:
period tests are approximate tests against the true margin, not a
claim that the candidate's raw periods are correct. When $e$ is below
the thresholds in that lemma, the recovered orbit count and rational
relations are the true ones and the matched lattice basis error is
$O(e)$. Choose one candidate body per recovered true orbit and repeat
these bodies under the reconstructed basis. Lemma~\ref{lem:smoothing}
gives $C^2$ support error $O(\sigma^{2/3})$, including the center error.
The basis and its inverse remain bounded. Consequently only a bounded
set of integer translates can enter any bounded fundamental region;
positive physical separation persists for small $\sigma$ and the
new presentation defines a valid periodic union. Its primitive free
area error is $O(\sigma)$ by \eqref{eq:area} and the area estimate.

Choose \eqref{eq:scalar-scale} with a small constant and decrease
$\nu_0$ to impose the rolling-radius, separation, patch-gap and rational
thresholds. This proves the geometric accuracy and transforms
\eqref{eq:raw-cost} into \eqref{eq:scalar-cost}. No additional
independence is needed for root or candidate selection, since every
conclusion holds on the one simultaneous deterministic-command event.
\end{proof}

The finite-presentation metric just used compares one matched body for
each true primitive component orbit, their support functions in the
common laboratory directions, and matched period bases. It does not
assert bounded Hausdorff distance between two entire infinite unions
whose periods differ slightly. Such a metric would encode a different
question. The exact rational relations are relations within the chosen
period-generating list; they are not exact Euclidean coordinate values.

\begin{remark}[Calibration and the meaning of the rate]\label{rem:resources}
A command implemented within total variation $\epsilon_0/2$ obeys the
bias allowance. With this kernel, spatial center error at most
$c_\kappa\epsilon_0\sigma$ is one sufficient contribution to that
allowance. The theorem supposes the two directions and the horizon
are calibrated; it does not prove a uniform angular-jitter bound near
grazing rays. Only their combined certified probability bias, when
available, may be charged to the same allowance. The output bit does
not certify input accuracy. Increasing spatial resolution therefore
still costs preparation precision, even though a fixed absolute accuracy
for each scalar mean suffices. This cost, the number of commands,
attempt count, bit arithmetic, and the infinite-class selection problem
are kept separate.
\end{remark}

\begin{corollary}[Unknown patch margin]\label{cor:pointwise}
For each fixed member of $\mathcal B$, a sequence of fresh finite
scalar experiments with $\sigma_j\downarrow0$ and summable simultaneous
failure probabilities gives convergent local geometry and eventually
correct primitive orbit count and rational period relations almost
surely. A threshold tending to zero more slowly than $\sigma_j$, such
as $\sqrt{\sigma_j}$, suffices for the patch comparisons. No uniform
finite stopping certificate over all of $\mathcal B$ is asserted.
\end{corollary}
\begin{proof}
Every relevant false translation in a fixed bounded patch has positive
defect, and there are finitely many such candidates. Let their minimum
be $\eta_{\mathcal O}>0$. The geometric errors are $O(\sigma_j)$,
so the proposed threshold eventually lies strictly between true and
false defects. All protected generators remain present, and the fixed
rational thresholds eventually hold. By the first Borel--Cantelli
lemma, only finitely many simultaneous mean events fail. Recompute
rather than permanently retain earlier classifications. The geometric
error estimate converges and the discrete decisions are eventually
correct. The unknown $\eta_{\mathcal O}$ does not supply an observable
finite index at which this has occurred.
\end{proof}

\begin{proposition}[Repeated motifs near a symmetry increase]
\label{prop:jump}
The patch-margin dependence cannot be removed from a uniform continuous
recovery assertion for the full period lattice and primitive orbit count
on the whole class $\mathcal B$.
\end{proposition}
\begin{proof}
Use the retained repeated-disk family \cite{Supplement}: disks of
radius $1/4$ centered at $(4k,4l)$ and $(4k+2+s,4l)$, where
$0\le s<1/8$ and $k,l\in\Z$. All displayed geometric and smoothness
bounds are uniform. At $s=0$ the full lattice is $2\Z\times4\Z$
and there is one component orbit; at $s>0$ it is $4\Z\times4\Z$
and there are two. Indeed a horizontal exchange translation preserves
$\{0,2+s\}$ modulo $4$ exactly when $2(2+s)$ is divisible by $4$.
The exchange candidate has central defect $2s$ for $s>0$. It can be
arbitrarily close to a period while not being one. The bodies converge
smoothly on every bounded patch but the exact discrete invariants jump.
This geometric example applies to any reconstruction derived from close
patch geometry. It does not contradict exact determination or the
pointwise conclusion above.
\end{proof}

\section{Calibration and randomized opposed launches}
\label{sec:calibration}
The scale-independent estimate in Proposition~\ref{prop:mean-stability}
concerns the propagation of mean error, not the implementation of an
input law. We first bound physical preparation errors without excluding
grazing commands. We then allow a calibrated distribution of directions
and durations rather than a single collimated displacement.

\subsection{A geometric bound for imperfect preparation}
For a compact convex body $C$ let
$S_a(C)=C+[-a,0]$. A segment $[q,q+a]$ hits $C$ precisely when
$q\in S_a(C)$. For a locally finite union of bodies this gives
\begin{equation}\label{eq:morph-bit}
 B_a(q)=\one_{\cup_C S_a(C)}(q)-\one_{\mathcal O}(q).
\end{equation}
The subtraction is essential: the first indicator counts a solid start
as a hit, whereas our recorded bit counts it as zero.

\begin{lemma}[Local convex boundary tubes]\label{lem:tube}
There is an absolute constant $c_0$ such that, for every planar compact
convex body $C$, disk $B_s(x)$ and $r>0$,
\[
 \area\{q\in B_s(x):\dist(q,\partial C)\le r\}
                 \le c_0(s+r)r.
\]
The bound also applies to the convex swept body $S_a(C)$, irrespective
of its straight sides or nonsmooth junctions.
\end{lemma}
\begin{proof}
Use signed distance, negative inside $C$, and integrate its level sets
between $-r$ and $r$ by the coarea formula. Its gradient has length one
almost everywhere away from the boundary and the interior medial set.
Positive levels are boundaries of outer parallel bodies; negative
levels are boundaries of convex erosions, until the erosion becomes
empty. The part of any such convex boundary in $B_s(x)$ has length at
most $2\pi s$: it is part of the boundary of its intersection with the
disk, whose perimeter is at most that of the disk by the planar Cauchy
perimeter formula. Coincident boundary pieces only reduce this bound.
Approximation by slightly larger disks handles tangencies and exceptional
levels. Integration gives $4\pi sr$; the displayed weaker bound also
covers all degenerate limiting cases.
\end{proof}

\begin{theorem}[Position, duration and angle calibration]\label{thm:calibration}
Assume the planar components are convex and mutually separated by
$d_0>0$. A nominal command has density $f_{x,\sigma}$, displacement
$a=tv$, $|v|=1$, $t\le b$, and $0<\sigma\le\sigma_0$.
Suppose an actual attempt can be coupled to this command so that
\[
 |q'-q|\le\ell,\qquad |t'-t|\le\tau,\qquad
 \angle(v',v)\le\alpha,
\]
and remains a unit-speed straight flight until its first collision.
Both nominal and actual protocols assign zero to solid starts and free
misses and count all attempts. Set
\begin{equation}\label{eq:calibration-size}
             r=\ell+\tau+b\alpha.
\end{equation}
If $r\le\sigma$, then
\begin{equation}\label{eq:calibration-bias}
 \left|\E B_{t'v'}(q')-P_a(f_{x,\sigma})\right|
                  \le C_{d_0,b,\sigma_0,\kappa}\,r/\sigma.
\end{equation}
The coupling errors may depend on $q$ and need not have mean zero.
An exceptional coupling event of probability $\zeta$ adds at most
$\zeta$ to the bound. The estimate is uniform in the nominal direction
and includes commands tangent to obstacles.
\end{theorem}
\begin{proof}
One has $|t'v'-tv|\le\tau+b\alpha$. The Hausdorff distance between
$S_{t'v'}(C)$ and $S_a(C)$ is at most this quantity. For convex bodies
$K,K'$ with $d_H(K,K')\le s$, the inclusions
$K_{-s}\subset K'\subset K+s\overline B$ follow directly from their
support inequalities. Thus a change of swept-body membership, including
the start displacement, can occur only when the nominal $q$ is within
$r$ of $\partial S_a(C)$. A change of solid-start membership is similarly
confined to the $r$-tube of $\partial C$. These containments hold for
every allowed actual displacement, even if its error depends on $q$.

Only a uniformly bounded number $N_0=N_0(d_0,b,\sigma_0)$ of components
can contribute such tubes in $B_\sigma(x)$. Indeed each contributing
component has a point in $B_{\sigma+b+r}(x)$; chosen points from distinct
components are separated by $d_0$, so disjoint radius-$d_0/2$ disks give
a uniform packing bound. Apply Lemma~\ref{lem:tube} to these bodies
and their sweeps in \eqref{eq:morph-bit}. Their exceptional area is at
most $C N_0(\sigma+r)r$. Since
$\|f_{x,\sigma}\|_\infty=\sigma^{-2}\|\kappa\|_\infty$ and
$r\le\sigma$, its probability is at most $Cr/\sigma$. Outside it
both bits agree. A failed coupling contributes at most its probability
because the bits belong to $[0,1]$.
\end{proof}

A systematic center error, angular error and timing error may therefore
all be charged explicitly to the mean-error budget. For fixed allowed
mean bias $b_0$, the sufficient condition is
$\ell+\tau+b\alpha\le c b_0\sigma$, not a scale-free apparatus
tolerance. No calibration estimate is obtained from the collision bits
themselves. Arbitrary curved pre-impact paths or an unknown distribution
outside this coupling model are not included.

\subsection{From one displacement to a stopped averaging operator}
Let $\rho$ be a known probability measure on displacements. In the
forward protocol draw $a\sim\rho$ and $q\sim f$ independently and
attempt $(q,a)$. In the reverse protocol draw the same law for $a$, draw
$q\sim f$, and attempt $(q+a,-a)$. These two protocols need not use
the same physical particle or the same random draw. Their prescribed
joint distributions must, however, have this translated relationship.
Only their mean bits are used. Write their means as $P_\rho^+(f)$ and
$P_\rho^-(f)$, and define
\[
 (T_\rho u)(x)=\int u(x+a)\rho(\dd a).
\]
Averaging Theorem~\ref{thm:balance} shows that their difference is
$\int f(T_\rho\chi-\chi)$; for the kernel commands it is
\begin{equation}\label{eq:random-balance}
       g_\sigma(x)=(T_\rho-I)u_\sigma(x).
\end{equation}
This is a coupled preparation design. Independently randomized reverse
starts without the translation by their sampled displacement do not
in general give \eqref{eq:random-balance}.

\begin{theorem}[Finite stopped inverse]\label{thm:stopped}
Let $0\le u\le1$, let $T$ be the transition operator of a Markov chain
$(X_j)$, and suppose that, for every starting point in the region in
question, the first visit to $\{u=0\}$ occurs by time $m$ almost surely.
Values of $u$ and $g=(T-I)u$ need only be defined on its $m$-step
reachable region. Put
\begin{equation}\label{eq:bellman}
 V_0=0,\qquad V_{n+1}=\max\{0,TV_n-g\},\qquad 0\le n<m.
\end{equation}
Then $V_m=u$. For a perturbation $\widehat g$ with uniform error
at most $\xi$, the same recursion, optionally clipped to $[0,1]$,
has error at most $m\xi$ at step $m$. Each additional uniform error
in evaluating $TV_n$ adds at most its magnitude to the final bound.

For $T=T_\rho$, the hypothesis holds for $u=\chi$ and $u=u_\sigma$
under the component bounds if, for some unit $v$ and $\gamma>0$,
\begin{equation}\label{eq:cone}
 a\cdot v\ge\gamma,
 \quad |a|\le b<d_0-2\sigma_0
 \quad(\rho\text{-almost every }a),
 \qquad m\gamma>D_0+2\sigma_0.
\end{equation}
Thus two averaged scalar-launch functionals with a fixed, possibly
continuous spread of durations and directions determine the local
obstacle set. Neither a realized impact position nor a free reference
point is supplied. The aperture enlargement is at most $mb+\sigma_0$.
\end{theorem}
\begin{proof}
Induction gives $0\le V_n\le u$, because $Tu-g=u$. At a point where
$u=0$, every $V_n$ is zero. At a point where $u>0$, writing
$e_n=u-V_n\ge0$ gives
$e_{n+1}=\min(u,Te_n)\le Te_n$. Consequently
\[
 e_m(x)\le\E_x\left[u(X_m)
                 \prod_{j=0}^{m-1}\one_{\{u(X_j)>0\}}\right]=0.
\]
The last equality uses the assumed visit by time $m$, including its
endpoint. Positivity of $T$ and nonexpansiveness of the maximum prove
$\|V_{n+1}-\widehat V_{n+1}\|_\infty
 \le\|V_n-\widehat V_n\|_\infty+\xi$.
Clipping cannot increase error relative to $V_n\in[0,1]$. This proves
the stability assertions.

Under \eqref{eq:cone}, a path whose first $m+1$ points have positive
$u_\sigma$ must stay in a single $\sigma$-enlarged component: successive
steps are shorter than the distance between distinct enlarged components.
Its endpoints differ by at least $m\gamma$ in direction $v$, exceeding
the diameter of that component. This is impossible. The indicator case
uses the same argument without enlargement. Approximate-identity
convergence after \eqref{eq:random-balance} recovers $\chi$.
\end{proof}

The recursion is the elementary finite-horizon obstacle recursion, here
with a physically supplied forcing $g$ and a geometrically bounded
zero-hitting time. For $\rho=\delta_a$ its expansion is exactly the
minimum of partial sums in \eqref{eq:min-inverse}. The finite stopping
condition cannot be discarded: for a symmetric nearest-neighbor walk
and $u=1$ on two adjacent sites and zero elsewhere, the two nonzero
values of $V_n$ are $1-2^{-n}$, not one at any finite $n$.
Theorem~\ref{thm:calibration} may be conditioned on the sampled nominal
displacement in either protocol. Uniformly bounded implementation errors
then perturb each averaged scalar mean by the same explicit bound.

\section{Constructive patch recovery and finite arithmetic}
\label{sec:constructive}
The selection in Theorem~\ref{thm:finite} is an existence regularization.
It is not an algorithm for searching an infinite-dimensional physical
class. We now give a different reconstruction from the same fixed-horizon
scalar commands. No minimization over physical candidates is used.

\begin{lemma}[Threshold clusters and convex hulls]\label{lem:hulls}
Let the complete convex components needed in a protected inner patch
have separation at least $d_0$, diameter at most $D_0$ and interior
rolling radius at least $r_*>0$. Let a square grid of covering radius
$\rho\le\sigma/8$ cover a surrounding square $U$, with the protective
collar of Lemma~\ref{lem:grid}. Assume
$\sigma<\min(r_*/4,d_0/20,1/4)$, and suppose that at every grid node
\begin{equation}\label{eq:grid-value-error}
                    |\widehat u(x)-u_\sigma(x)|\le1/8.
\end{equation}
Retain nodes with $\widehat u(x)\ge1/2$, join pairs at Euclidean distance
at most $4\sigma$, and take the convex hull of each graph component.
For every physical body in the protected patch there is exactly one
such cluster, and its hull $H_C$ satisfies
\begin{equation}\label{eq:hull-error}
                         d_H(H_C,C)\le2\sigma+\rho<3\sigma.
\end{equation}
No cluster joins two physical bodies. Outer fragments are excluded by
using only hulls at distance more than $D_0+1$ from $\partial U$;
all bodies needed on the smaller protected patch pass this cutoff.
\end{lemma}
\begin{proof}
A retained node is within $\sigma$ of a unique component, since otherwise
$u_\sigma=0$ and it cannot pass the threshold. Every grid node of
$C_{-\sigma}$ is retained, since its occupation average is one.
As in Lemma~\ref{lem:grid}, the rolling condition gives
$C=C_{-2\sigma}+2\sigma\overline B$.

For a retained node $x$ assigned to $C$, choose $y\in C_{-2\sigma}$
with $|x-y|\le3\sigma$, and a nearest grid node $z$. Then
$|z-y|\le\rho$ and $z\in C_{-(2\sigma-\rho)}\subset C_{-\sigma}$,
so $z$ is retained and $|x-z|<4\sigma$. Any two such deeper nodes can
be joined by a straight segment in the convex set
$C_{-(2\sigma-\rho)}$. Partition the segment into lengths at most
$\rho$ and replace the partition points by nearest grid nodes. The
new nodes belong to $C_{-(2\sigma-2\rho)}\subset C_{-\sigma}$ and
successive distances are at most $3\rho<4\sigma$. Thus the whole
cluster assigned to $C$ is connected. Nodes assigned to different
bodies are at distance at least $d_0-2\sigma>4\sigma$, so no edge
joins them.

Every retained node assigned to $C$ belongs to $C+\sigma\overline B$;
convexity gives the same inclusion for its hull. Conversely, each point
of $C$ is within $2\sigma$ of $C_{-2\sigma}$ and then within $\rho$
of a retained grid node. This proves \eqref{eq:hull-error}. The collar
ensures that all nodes and connecting paths used for an inner component
are observed. A component meeting $\partial U$ can contribute retained
nodes only within $D_0+\sigma$ of that boundary, and so cannot supply a
hull passing the stated cutoff. No unobserved completion of an outer
fragment is used.
\end{proof}

\begin{proposition}[Smooth reconstruction from a polygon]\label{prop:hull-smooth}
Suppose $C$ has a bounded centered $C^6$ support function, its center
lies in a bounded patch, and $d_H(H,C)\le e$ for a convex polygon $H$.
The polygon need not satisfy a smoothness prior. With the compensated
kernel $J$ of Lemma~\ref{lem:smoothing}, set $h=e^{1/6}$. Then
\begin{equation}\label{eq:polygon-smoothing}
 \|J_h*p_H-p_C\|_{C^2}\le C(eh^{-2}+h^4)\le Ce^{2/3},
\end{equation}
where here $p_H,p_C$ are the uncentered support functions in the same
laboratory frame. For sufficiently small $e$ this is the support
function of a smooth strictly convex body. Certified quadrature to
$C^2$ error $O(e^{2/3})$ gives the same assertion.
\end{proposition}
\begin{proof}
The polygon has $\|p_H-p_C\|_\infty\le e$. Differentiating the kernel,
not the polygon, bounds $J_h*(p_H-p_C)$ in $C^2$ by $Ceh^{-2}$.
Apply the moment cancellation estimate only to the smooth true $p_C$,
which gives the $Ch^4$ term. This does not assume a fictitious $C^6$
bound for the reconstructed polygon. The true radius of curvature
$p_C+p_C''$ has a positive lower bound. The same is therefore true
of the reconstruction when its $C^2$ error is small.

For specificity, use a periodic uniform quadrature mesh of width at
most $e^2$. The integrands $J_h^{(j)}(\theta-s)p_H(s)$, $j\le2$,
have uniformly bounded Lipschitz constants $Ch^{-4}$, since the bounded
polygon has a uniformly Lipschitz support function. The quadrature
error through two derivatives is $Ce^2h^{-4}=Ce^{4/3}$, uniformly in
$\theta$. The resulting finite sum of translates of a fixed smooth
kernel is itself smooth; its curvature radius remains positive. This
also supplies a finite representation of the reconstruction.
\end{proof}

\begin{theorem}[Constructive calibrated scalar reconstruction]
\label{thm:constructive}
Fix every numerical bound of $\mathcal B_\eta$, including a known
positive patch margin $\eta$. There are $C,c,\nu_0>0$ such that for
$0<\nu<\nu_0$ the following deterministic command design and finite
reconstruction have confidence at least $1-\delta$, $0<\delta<1/4$.
Use the two nominal directions and nominal duration in \eqref{eq:design},
localization $\sigma=c\nu^{3/2}$, and $O(\sigma^{-2})$ spatial
commands in the fixed aperture. Each independent attempt records one
bit; a solid start and a free miss both give zero and both remain in
the denominator. Allow actual position, duration and angular errors
satisfying the coupling model of Theorem~\ref{thm:calibration}, with
\begin{equation}\label{eq:constructive-calibration}
                    \ell+\tau+t\alpha\le c\sigma.
\end{equation}
With at most
\begin{equation}\label{eq:constructive-attempts}
                  C\nu^{-3}\log\frac{C}{\nu\delta}
\end{equation}
attempts, the reconstruction gives a smooth finite presentation with
matched $C^2$ support error, period-basis error and free-area error at
most $C\nu$. It recovers the true primitive orbit count and exact
rational relations of the true generating list in a true independent-pair
basis. Euclidean generator coordinates are estimated, not exact.

For fixed priors, $O(\sigma^{-4})=O(\nu^{-6})$ elementary arithmetic,
comparison and fixed-kernel evaluations suffice after the scalar means
are supplied. This conservative arithmetic bound is not an optimized
bit-complexity bound, a bound on apparatus motion, or a statistical
minimax assertion. No infinite-dimensional candidate selection occurs
in this reconstruction. A known positive $\eta$ is required for its
uniform finite discrete decisions.
\end{theorem}
\begin{proof}
Take a grid on $Q_{R_0}$ with covering radius at most
$\sigma/(64K)$, and command all the forward chain centers needed for
\eqref{eq:smoothed-min}. There are $N=O(\sigma^{-2})$ grid nodes and
at most $2mN$ scalar means. Estimate each actual mean to accuracy
$1/(128m)$. Decrease the constant in \eqref{eq:constructive-calibration}
so that the calibration bias is at most $1/(128m)$ as well. The mean
error relative to the ideal protocol is at most $1/(64m)$; hence the
finite-chain formula gives occupation error at most $1/32$. Reserve
another $1/32$ for numerical arithmetic. This is smaller than the
allowance in \eqref{eq:grid-value-error}. Hoeffding's inequality and a
union bound give this simultaneous event with
$C\log(CN/\delta)$ independent attempts per command. All commands
are fixed before sampling; no selection of an onset or free point is
required. The resulting cost is \eqref{eq:constructive-attempts}.

Apply Lemma~\ref{lem:hulls}. The aperture margins in
\eqref{eq:aperture} ensure that every body with center in $Q_{10B}$
is complete and survives the exterior cutoff. All resulting hulls
have Hausdorff error $e\le3\sigma$. Compute their Steiner centers and
support suprema with certified numerical enclosures of error $O(\sigma)$.
The geometric tests and rational search in Lemma~\ref{lem:locking}
apply to these convex hulls: that lemma requires Hausdorff control,
not smoothness of the approximating body or admissibility of an entire
candidate table. For $\sigma$ below its $\eta$-dependent thresholds,
the orbit classifications and rational relations are correct, and the
matched period basis has error $O(\sigma)$.

Use one hull per recovered primitive orbit. Proposition~\ref{prop:hull-smooth}
turns it into a genuine smooth strictly convex body with $C^2$ error
$O(\sigma^{2/3})$. Repeat these bodies under the computed basis.
Its inverse remains bounded, so only a bounded number of integer
translates can enter any bounded fundamental patch; the true separation
therefore persists for small $\sigma$. This gives a physical periodic
presentation in the finite-presentation metric of Section~\ref{sec:finite}.
The polygonal area estimates and lattice covolume have error $O(\sigma)$;
the area of the final smoothed presentation has error
$O(\sigma^{2/3})$, also at most $C\nu$. These are different area
estimates and are not conflated.

For completeness we account for a finite implementation. The chain
sums cost $O(mN)$. Testing all pairs of retained grid nodes costs
$O(N^2)$, and connected components and monotone-chain convex hulls add
at most $O(N^2+N\log N)$. A hull support value is the maximum of its
vertex scalar products. Periodic angular meshes of width $O(\sigma^2)$
therefore compute all center integrals and support suprema to the
reserved $O(\sigma)$ precision in at most $O(N\sigma^{-2})$
operations. There are only a prior-bounded number of complete bodies
and period candidates. The bounded-denominator search and Hermite
calculation have prior-bounded size. The same angular mesh supplies
the finite smooth-kernel representation of
Proposition~\ref{prop:hull-smooth}, again in $O(N\sigma^{-2})$
operations. Derivative and curvature tests can be enclosed to the
reserved positive margins. Since $N=O(\sigma^{-2})$, the total is
$O(\sigma^{-4})$. Fixed elementary functions and the chosen kernel
are assumed available with certified evaluation; the count is in this
explicit arithmetic model, not a claim about unrestricted real-number
oracles in a Turing model. This completes the proof.
\end{proof}

The new finite procedure and the retained candidate-based procedure
have different computational claims. Both retain the same conditional
statistical order and spatial input resolution. The former removes the
unimplemented search step, not the geometric priors or calibration cost.
With unknown $\eta$, recomputation at successively finer scales gives
the pointwise eventual conclusion of Corollary~\ref{cor:pointwise},
not an observable uniform stopping certificate. The symmetry-increase
example of Proposition~\ref{prop:jump} remains applicable.

\section{Information categories, local periodicity and retained results}
\label{sec:comparison}
\subsection{What is observed and what is reconstructed}
There are three different collision experiments in the present submission
package. The spatial first-impact experiment of Supplement P launches
uniformly in a fixed laboratory square with independent uniform directions
and records a localized first-impact position. In exact data, the support
of its position law is the relevant physical boundary union. Its protected
launch collars give positive mass to every boundary arc. The exact
finite-patch inverse in that experiment is therefore informationally
direct; it is not boundary recovery from hidden travel-time or spectral
data.

The experiment of this primary article records no collision coordinates.
It instead allows localized commanded launch densities and two opposed
collimated directions. The measured imbalance is a finite difference
of occupation, not the occupation function itself. The cancellation in
Theorem~\ref{thm:balance} removes paths which enter and leave a body
between two free endpoints. The finite-chain zero condition restores
the additive information lost in taking that difference. The same
nonlinear inverse works for mollified occupation without differentiating
noisy probability functions. The theorem is thus a scalar-output
alternative within the collision-law programme, with an explicitly
richer preparation contract than a single uniform launch law. Neither
experiment is identified with the other by an unsupported information
ordering.

Earlier intrinsic return-law results concern different selected
itineraries, endpoint arclength laws and moment compressions. Their
count-only finite fibers, formal fibers and smooth Taylor agreements
remain distinct from equality of exact analytic count germs. A family
of spatially commanded binary collision probabilities is not the
previously studied uniformly prepared count-only datum. The present
argument makes no inference about that analytic count-germ problem,
passive trajectories, marked-length spectra or spectral rigidity.

\subsection{Finite differences and local periodicity}
The finite-difference operator in \eqref{eq:balance} is elementary.
The point is the physical identity realizing it without output positions,
and the bounded zero-witness inverse, rather than a new general theorem
about difference operators. Covariogram formulas also link translations
to occupation statistics: for a measurable finite-volume set, the
covariogram integrates the product of two translated indicators.
Galerne~\cite{Galerne} relates its directional derivatives to perimeter
and directional variation. Here spatially specified test densities
retain the location of a \emph{signed difference}; the datum is not a
translation-invariant autocorrelation. No covariogram uniqueness or
priority conclusion follows from the comparison.

There is also a distinction between recognizing a period and proving
that a set must be periodic. Lagarias and Pleasants~\cite{LP} study
patch-counting complexity and crystallinity of Delone sets. Dolbilin,
Garber, Schulte and Senechal~\cite{DGSS} study regularity radii: local
cluster equivalence can impose global regularity, with quantitative
bounds in the classes they treat. Herva and Kari~\cite{HK} develop
local-complexity and annihilator methods for Delone configurations,
including periodic decomposition and criteria forcing periodicity.
These are local-to-global structural questions with their own hypotheses.

In contrast, \eqref{eq:period-presentation} already assumes a bounded
periodic presentation. Its bound ensures a known finite central patch
contains representatives of every orbit and a protected list containing
generators of the full period group. Our period test does not infer
crystallinity for an arbitrary Delone set from observing a large patch.
The new scalar balance need not vanish and is not an annihilator
hypothesis. Its inversion first recovers the physical patch; only then
is the previously proved finite translation criterion used. This
separation places the contribution without suggesting that bounded
period recognition supersedes general local-periodicity theory.

\subsection{Two finite experiments and their different resources}
On the repeated-motif patch-margin class, the retained spatial first-hit
theorem gives a sufficient attempt bound
$C_\eta\nu^{-3/4}\log(C_\eta/(\nu\delta))$, with output localization
$O(\nu^{3/2})$, for its uniform-launch boundary-coverage design. The
scalar theorem here gives
$C_\eta\nu^{-3}\log(C_\eta/(\nu\delta))$, with input localization
$O(\nu^{3/2})$ and no position output, using a two-dimensional grid of
commands and a fixed precision per scalar mean. Both constants and
accuracy thresholds depend on the same kind of nonperiod patch margin,
in addition to their respective apparatus and geometric bounds. These
are sufficient upper bounds for different experiments, not competing
minimax rates. The former's boundary coverage and the latter's scalar
spatial grid count different operations.

The boundary smoothing estimate is a compensated-convolution argument,
and the probability estimate is an elementary bounded-variable inequality.
Neither exponent is offered as a separate assertion of global optimality.
The exact reversal identity, finite-chain inverse and the explicit
preparation/output trade are the additions; the geometric period lemma
and arithmetic classification retain their preceding provenance.

\subsection{Submission architecture and source status}
Supplement P is the entire reviewed A2 v28 paper and its nested source
history, supplied without changing any mathematical input. It includes
the repeated-motif and translation-separated first-impact theorems,
record-local collision clouds, intrinsic return-law inverses, completion,
sequential certificates, moment theory and the relative-law supplementary
volumes. Those results are neither deleted nor asserted anew by the
primary article. The present text includes the period proof it uses
and has no dependency on the older return-law chapters.

The phrase ``already published'' in the preserved v28 README refers to
a repository deposit and must not be read as journal publication of v27.
The corrected status in this submission is \emph{previously deposited
and source-qualified author revision}. A prior exact-source build
qualifies its own frozen source only; the new primary and retained
volumes need the current qualification receipt. These source statements
do not replace independent proof review or a publication record.

\section{Capacity functionals and active probing}
\label{sec:hit-comparison}
The relation between intersection probabilities and geometry predates
the collision experiment studied here. Matheron~\cite{Matheron} formulates
stereological and morphological information using probabilities that a
structuring figure hits a random set or is contained in it. The capacity
functional $T_X(K)=\Prob\{X\cap K\ne\varnothing\}$ is a central
object in the theory of random closed sets; see Molchanov~\cite{Molchanov},
Chapter~1. The following identity places our observation in that
framework without identifying its endpoint convention with an ordinary
hit test.

\begin{proposition}[The endpoint-excluded capacity datum]
\label{prop:capacity}
Let $Q$ have the commanded density $f$, let $\mathcal O$ be a fixed
closed set, and let $X=\mathcal O-Q$. Then
\begin{equation}\label{eq:capacity-bit}
 P_a(f)=T_X([0,a])-T_X(\{0\}).
\end{equation}
For $X_a=\mathcal O-(Q+a)$ the opposed translated experiment satisfies
\[
 T_{X_a}([-a,0])=T_X([0,a]),\qquad
 P_a(f)-P_{-a}(f_a)=T_{X_a}(\{0\})-T_X(\{0\}).
\]
The equality holds without a convexity or periodicity hypothesis.
\end{proposition}
\begin{proof}
The event $0\in X$ is contained in the event $X\cap[0,a]\ne\varnothing$.
Their difference is precisely a free start with a segment hit. In the
opposed experiment the tested segment has the same physical image, while
its starting endpoint is $Q+a$. Subtraction gives the second formula.
No conditional probability is introduced: solid-start attempts remain
among all draws of $Q$.
\end{proof}

Thus the affine family of capacity probabilities itself is not a new
notion, and the morphology of $C+[-a,0]$ is not a new set operation.
The endpoint exclusion cancels the common segment-hit term and exposes
a signed spatial occupation difference. The bounded zero witness makes
that difference invertible on the specified physical class. The stopped
averaging inverse extends this cancellation to coupled displacement
mixtures. Neither averaging nor a finite-horizon obstacle recursion is
claimed as a new abstract probability principle; the statements specify
which experimentally accessible forcing and which geometric stopping
condition give exact recovery and a finite error bound.

Active probing supplies a second nearby comparison, with different
response alphabets. Edelsbrunner and Skiena~\cite{ES} reconstruct convex
polygons by X-ray probes reporting the \emph{length} of intersection
with a line; their determination bound is $5n+O(1)$ probes for an
$n$-vertex polygon, and their verification bounds concern a given
polygon. A real intersection length is not our Bernoulli collision
output. Their finite-dimensional polygonal target and adaptive exact
line probes are also different from the smooth-class probability
estimation and localized spatial inputs used here. One cannot compare
those probe counts with \eqref{eq:constructive-attempts} while ignoring
both the unknown class and the information returned by a probe.

Bose, De Carufel, Shaikhet and Smid~\cite{BCSS} give a reconstruction
algorithm from wedge probes. Their Theorem~4.1 assumes
$0<\omega\le\pi/2$ and every polygon angle greater than $\omega$,
and obtains at most $2n-2$ probes. A probe reports the wedge position,
its rays and boundary contact points. Our experiment supplies no contact
point and returns one bit per attempt, but pays for a growing grid of
spatially localized commands and repetitions. Consequently neither
observation is asserted to dominate the other in an information order.
These comparisons concern the cited probe models and results, not an
exhaustive priority search in geometric tomography.

The contribution of the present revision is accordingly precise. It
couples the scalar occupation inverse to a stopped randomized-preparation
version, proves a grazing-uniform apparatus-error bound, and replaces
one infinite-class existence selection by a proved finite geometric
procedure. The bounded-period and patch-margin hypotheses remain visible
in the global conclusion. No statistical lower bound or optimality,
passive rigidity theorem, discovery of periodicity without a prior, or
exact analytic uniform-count theorem is inferred from these additions.

\subsection*{Current submission map}
The five core chapters of the reviewed scalar article remain active and
unchanged. The complete reviewed article, including its validation files
and all nested earlier volumes, is also supplied at \texttt{retained/v29}.
In that retained terminology, Supplement P denotes v28; its current path
is \texttt{retained/v29/retained/v28}. The present primary adds the
calibration, stopped inverse, constructive recovery and capacity comparison
as self-contained arguments. Earlier existence-only statements describe
their original algorithms and are not relabelled as computational bounds
for this revision. The source map is a preservation record, not a claim
that a former build receipt validates the new proofs.

\section{Reciprocal inversion without a common drift direction}
\label{sec:mean-exit}
The finite stopping hypothesis of Theorem~\ref{thm:stopped} is sufficient
but is not necessary for identification. A centered random displacement
can remain inside one component for arbitrarily many steps. Nevertheless,
its mean exit time is uniformly bounded by the component diameter and
its second moment. We use that bound twice: to recover the occupation
without specifying its zero set, and to compare two unknown occupations
with a constant independent of the iteration depth.

\subsection{The observed forcing and the computational walk}
Let $\rho$ be a prescribed probability measure on $\R^d$. The forward
protocol draws $q\sim f$ and $a\sim\rho$ independently and attempts the
straight segment from $q$ to $q+a$. The reverse protocol draws the same
joint law for $(q,a)$ but attempts the segment from $q+a$ to $q$.
Each protocol records only its hit bit. A solid start and a free miss
are both zero; all attempts are counted. No displacement-resolved
subpopulation mean is required. The direction and duration randomizer
is part of the input controller, not an additional geometric output.

With $\chi=\one_{\mathcal O}$ and
$T_\rho w(x)=\int w(x+a)\rho(\dd a)$, the pointwise reversal balance
and Fubini give
\begin{equation}\label{eq:reciprocal-forcing}
 P_\rho^+(f)-P_\rho^-(f)=\int f(T_\rho-I)\chi.
\end{equation}
For the translated even kernel commands of Section~\ref{sec:acquisition},
translation invariance commutes with convolution. Thus the observed
balance function is $g_\sigma=(T_\rho-I)u_\sigma$, where
$u_\sigma=\kappa_\sigma*\chi$. The corresponding random walk is a
\emph{computational} process. Its successive locations are not a
multi-collision trajectory and its exit from an unknown component is
not an observed stopping signal.

In this section impose, for known $b,v_*>0$,
\begin{equation}\label{eq:centered-law}
 |a|\le b\quad\rho\text{-a.s.},\qquad
 \int a\,\rho(\dd a)=0,\qquad
 s_2:=\int|a|^2\rho(\dd a)\ge v_*.
\end{equation}
There need not be a direction with positive drift. Neither symmetry,
a density for $\rho$, nor a positive definite covariance matrix is
required. The full angular circle, or two equally weighted opposite
steps, satisfies these conditions when the nonzero step length is fixed.
The reverse translation in \eqref{eq:reciprocal-forcing} remains essential;
independently choosing a reverse position with the same marginal
angular law is a different experiment.

\begin{lemma}[A uniform mean exit bound]\label{lem:mean-exit}
Let $0\le u\le1$ be measurable. Suppose its positive set is contained
in a union of sets of diameter at most $D$, separated from one another
by a distance strictly greater than $b$. Let $X_n=x+a_1+\cdots+a_n$,
where the increments are independent with law \eqref{eq:centered-law},
and put $\tau_u=\inf\{n\ge0:u(X_n)=0\}$. Then
\begin{equation}\label{eq:exit-bound}
 \sup_x\E_x\tau_u\le H:=\frac{(D+b)^2}{v_*},\qquad
 \sup_x\Prob_x(\tau_u>n)\le2^{-\lfloor n/L\rfloor},
 \quad L=\max(1,\lceil2H\rceil).
\end{equation}
No regularity of the positive set is used.
\end{lemma}
\begin{proof}
If $u(x)=0$, the exit time is zero. Otherwise, until that time all
successive locations with positive $u$ lie in the same containing set:
a step of length at most $b$ cannot jump to a different one. For
$S=\tau_u\wedge n$ we consequently have $|X_S-x|\le D+b$.
Independence and zero mean make
$|X_n-x|^2-ns_2$ a martingale. Applying the bounded stopping identity
at $S$ gives
\[
 s_2\E_x S=\E_x|X_S-x|^2\le(D+b)^2.
\]
Monotone convergence proves the expectation estimate, including
almost-sure finiteness of $\tau_u$. Markov's inequality gives
$\Prob_x(\tau_u>L)\le H/L\le1/2$. At time $L$, on survival,
the same expectation estimate applies from the current location.
Iteration of the Markov property gives the stated geometric tail.
This argument does not presume optional stopping at an unbounded time;
only bounded stopping is used before taking the limit.
\end{proof}

\begin{theorem}[Mean-exit inverse and physical comparison]
\label{thm:mean-exit-inverse}
Under Lemma~\ref{lem:mean-exit}, write $g=(T_\rho-I)u$ and define
\begin{equation}\label{eq:mean-exit-recursion}
 V_0=0,\qquad V_{n+1}=\max\{0,T_\rho V_n-g\}.
\end{equation}
Then $0\le V_n\le V_{n+1}\le u$ and
\begin{equation}\label{eq:mean-exit-rate}
 \|u-V_n\|_\infty\le2^{-\lfloor n/L\rfloor}.
\end{equation}
If $\widehat g$ has uniform error at most $\xi$, and each transition
update has additional uniform error at most $\zeta$, clipping the
perturbed iterates to $[0,1]$ gives
\begin{equation}\label{eq:mean-exit-noise}
 \|u-\widehat V_n\|_\infty
       \le2^{-\lfloor n/L\rfloor}+n(\xi+\zeta).
\end{equation}
In particular, increasing the iteration depth without regard to the
error is not the asserted regularization.

There is a stronger comparison between physical candidates. Suppose
$u,\widetilde u\in[0,1]$ both satisfy the same containing-set bounds,
with no prescribed agreement of their zero sets. If
$g=(T_\rho-I)u$ and $\widetilde g=(T_\rho-I)\widetilde u$, then
\begin{equation}\label{eq:unknown-support-stability}
 \|u-\widetilde u\|_\infty
                  \le H\|g-\widetilde g\|_\infty.
\end{equation}
Thus $g$ determines $u$ in this class. This estimate is independent of
both a convolution scale and the iteration depth. It compares admissible
occupations, not arbitrary noisy outputs of the recursion.
\end{theorem}
\begin{proof}
Positivity of $T_\rho$ and $T_\rho u-g=u$ give the monotone bounds.
Every iterate vanishes where $u=0$. For $e_n=u-V_n$, at a point with
$u>0$ the recursion gives $e_{n+1}=\min(u,T_\rho e_n)\le T_\rho e_n$.
Iterating until the first zero yields
\[
 0\le e_n(x)\le\E_x\{u(X_n)\one_{\{\tau_u>n\}}\}.
\]
Lemma~\ref{lem:mean-exit} proves \eqref{eq:mean-exit-rate}.
The maximum, clipping and a Markov transition are nonexpansive in
supremum norm. Induction proves \eqref{eq:mean-exit-noise}.

For the comparison put $w=u-\widetilde u$ and
$\epsilon=\|g-\widetilde g\|_\infty$. Bounded stopping in the
martingale for $w$ gives, with $S=\tau_u\wedge n$,
\[
 w(x)=\E_x w(X_S)
       -\E_x\sum_{j<S}(g-\widetilde g)(X_j).
\]
At $\tau_u$ the terminal difference is $-\widetilde u\le0$; on
survival its positive part is at most one. Hence
$w(x)\le\Prob_x(\tau_u>n)+\epsilon\E_x(\tau_u\wedge n)$.
Letting $n$ increase gives $w(x)\le H\epsilon$. Interchanging the
two functions and stopping at $\tau_{\widetilde u}$ proves the reverse
bound. This uses neither an observed free reference point nor knowledge
of either stopping set in the reconstruction.
\end{proof}

\begin{corollary}[Every nontrivial bounded displacement law]
\label{cor:noncentered-exit}
The same identification and fixed-aperture comparison hold for any
prescribed law supported in $|a|\le b$ other than $\delta_0$, provided
the containing sets have separation greater than $b$. More explicitly,
suppose for some unit vector $v$ and numbers $\gamma,p>0$ that
$\rho\{a:a\cdot v\ge\gamma\}\ge p$. Choose an integer $m$ with
$m\gamma>D$. Then
\[
 \Prob_x(\tau_u>n)\le(1-p^m)^{\lfloor n/m\rfloor},\qquad
 \E_x\tau_u\le m/p^m.
\]
In Theorem~\ref{thm:mean-exit-inverse} and
Proposition~\ref{prop:fixed-exit-aperture}, replace the survival bound
and $H$ by these quantities. The centered-law bound
\eqref{eq:exit-bound} is usually substantially better; the constants
here can be large when the escape cone has small probability.
\end{corollary}
\begin{proof}
Conditional on any surviving starting position, the next $m$
independent increments all belong to the displayed cone with probability
at least $p^m$. If the occupation stayed positive throughout those
steps, they would lie in one containing set and their net projection
would exceed its diameter. Hence this event forces a visit to zero.
The Markov property at block endpoints gives the tail. Summing the tail
in blocks of length $m$ gives the expectation bound. The preceding
Bellman and stopped-comparison proofs require only these two bounds.
Finally, any law different from $\delta_0$ has a nonzero point in its
support; a sufficiently small ball about that point lies in a cone
$a\cdot v\ge\gamma>0$ and has positive probability. Thus the stated
condition does not require an almost-sure common drift direction.
\end{proof}

For mollified occupations under the physical component bounds, take
\begin{equation}\label{eq:uniform-exit-prior}
 D=D_0+2\sigma_0,\qquad b<d_0-2\sigma_0,
 \qquad 0<\sigma\le\sigma_0.
\end{equation}
Their positive sets lie in the $\sigma$-enlarged components, which
satisfy these diameter and separation bounds. Thus the same $H,L$
work at every scale. Equality of the exact two functionals on smooth
spatial commands determines every $g_\sigma$, hence every $u_\sigma$.
Approximate-identity convergence determines $\chi$ almost everywhere.
For the closed convex physical components this determines the sets
as closures of their interiors. Analyticity and periodicity are not
needed in this local inference.

\subsection{Locality, stability of the prescribed law, and sharpness}
An infinite iteration must not silently turn a local experiment into
an unbounded aperture experiment. The following localization resolves
that issue.

\begin{proposition}[A fixed aperture for the mean-exit inverse]
\label{prop:fixed-exit-aperture}
Let $U$ be a bounded target set and let a bounded observation set $W$
contain $U+\overline B_{D+b+\varepsilon_0}$ for some
$\varepsilon_0>0$. Define $T^Wv=T_\rho(\one_Wv)$ and perform
\eqref{eq:mean-exit-recursion} on $W$, setting every iterate to zero
outside $W$. Then \eqref{eq:mean-exit-rate} and
\eqref{eq:mean-exit-noise} hold on $U$, using only forcing errors on
$W$. The comparison \eqref{eq:unknown-support-stability} on $U$ likewise
requires forcing data only on this fixed enlarged set.
\end{proposition}
\begin{proof}
Globally the zero extension gives $T^Wu\le T_\rho u$, so induction
still gives $0\le V_n\le u$ and zero values on $\{u=0\}$.
For a start $x\in U$ with $u(x)>0$, the entire containing set of $x$
and every possible exit location lie in $W$. Along the walk up to
its first zero, $T^W$ therefore introduces no boundary discrepancy.
The killed error and stopped comparison proofs apply verbatim along
this walk. Errors elsewhere in $W$ can propagate through at most $n$
updates, giving the same bound as before. The argument on $U$ does
not require the unknown occupation to vanish near $\partial W$.
\end{proof}

Define total variation by
$\|\rho-\rho'\|_{\rm TV}=\sup_{0\le h\le1}|\int h\,\dd(\rho-\rho')|$.
For a physical $u\in[0,1]$,
\begin{equation}\label{eq:operator-calibration}
 \|(T_\rho-T_{\rho'})u\|_\infty
                    \le\|\rho-\rho'\|_{\rm TV}.
\end{equation}
Consequently an independently certified error in a reciprocal randomized
law can be charged directly to the forcing error. A baseline law
satisfying \eqref{eq:centered-law} suffices for this perturbation
comparison; the small error need not itself have mean zero. The actual
forward and reverse preparations must still be reciprocal. Uncoupled
position, timing and angle errors are instead charged by
Theorem~\ref{thm:calibration}, with its $\sigma^{-1}$ cost.
We do not infer either calibration guarantee from the observed bits.

The expectation factor has a familiar potential-theoretic meaning. If
$Q$ is a finite killed transition matrix with finite mean exit times,
then $G=(I-Q)^{-1}=\sum_{n\ge0}Q^n$ has nonnegative entries and
\begin{equation}\label{eq:green-norm}
             \|G\|_{\infty\to\infty}=\max_x\E_x\tau.
\end{equation}
Indeed its row sums are the sums of survival probabilities. The constant
is attained by forcing identically one on the killed state set. This is
a sharp Dirichlet resolvent norm, not a minimax lower bound for the
physical collision experiment with unknown obstacles. Green functions
and stopped martingales are classical tools; see \cite{LawlerLimic}.
Our use of them concerns a measured forcing and an unknown zero set,
with the containing sets needed only for bounds and not given as data.

For the symmetric nearest-neighbor walk on the line and $u=1$ at two
adjacent sites and zero elsewhere, \eqref{eq:mean-exit-recursion} gives
$V_n=1-2^{-n}$ at those sites. The mean exit time is two, although no
finite almost-sure stopping bound exists. This example belongs to the
new theorem and verifies why it is a genuine extension of
Theorem~\ref{thm:stopped}, not a claim that its earlier finite stopping
obstruction has disappeared. If the second moment vanishes, constants
on a stationary orbit cannot be recovered; the positive moment in
\eqref{eq:centered-law} has a different role from a common positive drift.

\section{A finite experiment with pooled reciprocal directions}
\label{sec:pooled-finite}
The mean-exit argument includes continuous angular distributions, but
that exact assertion alone does not give a finite quadrature algorithm.
Here we specify a finite-support law for which every transition is an
exact grid operation. Direction labels need not be retained with the
output, and no infinite-dimensional physical candidate search occurs.

Fix $t=\min(1,d_0/4)$ and the scale bound $\sigma_0$ in
\eqref{eq:design}. Use the four-direction mixture
\begin{equation}\label{eq:compass-law}
 \rho_\square=\tfrac14(\delta_{te_1}+\delta_{-te_1}
                        +\delta_{te_2}+\delta_{-te_2}).
\end{equation}
It has zero mean, second moment $t^2$ and no common positive drift
projection. Set
\begin{equation}\label{eq:pooled-design}
 H_*=(D_0+2\sigma_0+t)^2/t^2,\quad
 L_*=\lceil2H_*\rceil,\quad n_*=6L_*,\quad
 M_\sigma=\lceil32t/\sigma\rceil,\quad h_\sigma=t/M_\sigma.
\end{equation}
The grid $h_\sigma\Z^2$ has covering radius at most $\sigma/8$,
and all four displacements are grid translations by $M_\sigma$ indices.
These choices depend only on the priors and requested scale. In
particular $n_*$ does not grow as $\sigma$ decreases.

\begin{theorem}[Constructive pooled-direction recovery]
\label{thm:pooled-recovery}
Fix all numerical bounds of $\mathcal B_\eta$, including its known
positive patch margin $\eta$. For sufficiently small $\nu>0$ take
$\sigma=c\nu^{3/2}$ and the two reciprocal protocols
\eqref{eq:reciprocal-forcing} with law \eqref{eq:compass-law}.
At each commanded grid center estimate only the two \emph{pooled}
mean bits. No direction-resolved probabilities or collision positions
are observed. Solid-start attempts and free misses are both zero and
both are counted in the denominator.

A fixed-aperture finite algorithm constructs a smooth periodic
presentation with matched $C^2$ support error, period-basis error and
free-area error at most $C\nu$, and identifies the primitive orbit
count and exact rational generator relations with confidence
$1-\delta$. There are $O(\nu^{-3})$ spatial centers and at most
\begin{equation}\label{eq:pooled-budget}
                  C\nu^{-3}\log\frac{C}{\nu\delta}
\end{equation}
attempts. Constants depend on the stated class, including $\eta$.
Each ideal mean may have independently certified implementation bias
at most $1/(512n_*)$. For example, the position/timing/angle coupling
of Theorem~\ref{thm:calibration} satisfies this budget when
$\ell+\tau+t\alpha\le c\sigma/n_*$. A reciprocal-law total-variation
error can instead be charged as in \eqref{eq:operator-calibration}.

After the means are supplied, $O(\nu^{-6})$ arithmetic, comparison and
fixed-kernel evaluations suffice in the model of
Theorem~\ref{thm:constructive}. This includes $O(n_*\nu^{-3})$
local transition updates, not an exponentially growing tree of
randomized paths. No statistical optimality or Turing bit-complexity
claim is part of this theorem. Exact rational relations do not make
the estimated Euclidean basis coordinates exact.
\end{theorem}
\begin{proof}
Let the target grid cover $U=Q_{R_0}$ with $R_0$ as in
\eqref{eq:aperture}. Command both pooled means on a surrounding square
$W$ containing $U+\overline B_{D_0+2\sigma_0+t+1}$.
The actual forward and translated reverse starts lie in a further
collar of width $t+\sigma_0$. This fixed aperture is independent of
$\nu$, and has $N=O(\sigma^{-2})$ grid nodes.

Estimate each actual mean to accuracy $1/(512n_*)$. Together with the
allowed bias its error relative to the ideal mean is at most
$1/(256n_*)$. Their difference $\widehat g$ therefore has error at
most $1/(128n_*)$. Starting from zero, apply the clipped finite-grid
recursion
\[
 \widehat V_{j+1}(x)=\left[
   \tfrac14\sum_{a\in\{\pm te_1,\pm te_2\}}
        \widehat V_j(x+a)-\widehat g(x)\right]_{[0,1]},
 \qquad j<n_*,
\]
using zero extension outside $W$. No wraparound boundary is used.
The four transitions land exactly on the grid. On that grid, the
true values $u_\sigma(x)$ have the same containing-set and mean-exit
bounds as in the continuum. Proposition~\ref{prop:fixed-exit-aperture}
therefore gives error on $U$ at most
\[
       2^{-6}+1/128=3/128.
\]
Another $1/128$ may be reserved for finite numerical updates. The
resulting error $1/32$ is strictly below the threshold allowance $1/8$
in Lemma~\ref{lem:hulls}. Apply that lemma, then
Proposition~\ref{prop:hull-smooth} and Lemma~\ref{lem:locking}.
These provide Hausdorff error $O(\sigma)$, smooth support error
$O(\sigma^{2/3})$, and the claimed discrete and period conclusions.
The same protected component cutoffs, known $\eta$ thresholds and
rational separation bounds are used, not new assumptions of unequal
shapes or absent individual symmetries. The final smoothed free-area
error is $O(\sigma^{2/3})$ as in Theorem~\ref{thm:constructive}.

Each pooled attempt remains a Bernoulli trial even though its input
has been randomized. Hoeffding's inequality and a union bound over the
$2N$ commanded means give the simultaneous accuracy event with
$Cn_*^2\log(CN/\delta)$ attempts per mean. Since $n_*$ is fixed by
the priors, this proves \eqref{eq:pooled-budget}. Independence refers
to fresh attempts, not to all counterfactual outcomes on one segment.

Every Bellman layer uses four lookups per node; only two layers need be
stored. It costs $O(n_*N)$ arithmetic operations. The finite cluster,
hull, support and period stages have the $O(N^2)$ bound proved in
Theorem~\ref{thm:constructive}. Thus the complete arithmetic order is
$O(\sigma^{-4})$. This counting does not conceal a continuous averaging
oracle: the prescribed law has four atoms and the grid is closed under
its shifts before the explicit aperture killing.
\end{proof}

\subsection{Interval implementation and the role of calibration}
The finite transition step admits exact rational enclosures. If
$g(x)\in[g^-(x),g^+(x)]$ on the observed grid, start $V_0^-=V_0^+=0$
and use
\begin{equation}\label{eq:interval-bellman}
 V_{j+1}^-=[T^W V_j^- -g^+]_{[0,1]},\qquad
 V_{j+1}^+=[T^W V_j^+ -g^-]_{[0,1]}.
\end{equation}
Monotonicity implies that the true exact iterate lies between these
bounds. On the protected grid the occupation therefore belongs to
\begin{equation}\label{eq:occupation-enclosure}
 [\,V_n^-,\ \min(1,V_n^++2^{-\lfloor n/L_*\rfloor})\,].
\end{equation}
Empirical proportions, rational confidence allowances and the four
weights give rational endpoints. Thus the local stage can avoid
uncontrolled floating-point rounding. These enclosures are conditional
on the simultaneous mean-accuracy event and the geometric and reciprocal
preparation hypotheses. They do not validate an apparatus or certify
those priors from the bit record. Certified kernel evaluation remains
a separate requirement of the later smooth geometric stage.

For noisy refinement beyond threshold recovery, a valid iteration depth
balances the two terms in \eqref{eq:mean-exit-noise}; depth tending to
infinity at fixed forcing error is not justified. The physical-candidate
bound \eqref{eq:unknown-support-stability} has no such depth factor,
but does not itself make arbitrary measured forcing physically
admissible. This distinction prevents the comparison theorem from being
used as an unjustified stopping certificate for noisy Bellman iterates.

Optimal stopping, killed Green operators and value iteration have a
well-developed theory. The finite-expected-termination notion in
\cite{Bertsekas} concerns stochastic control problems with specified
costs and termination; our unknown is an occupation whose zero set is
not provided. The collision experiment supplies its signed forcing,
and the component bounds supply a termination estimate for the proof.
We use elementary parts of that machinery, not a new general Bellman
theorem. The new observation conclusion is the recovery from two pooled
reciprocal spatial-input functionals without a common directional cone,
and its fixed-aperture, finite-grid realization. The capacity and active
probing comparisons of Section~\ref{sec:hit-comparison} continue to
apply. In particular this is not a uniformly prepared count-law or
passive spectral inverse.

\subsection*{Acknowledgment and responsibility}
AI-assisted analysis contributed to the mean-exit comparison, fixed-aperture
localization, pooled-direction reconstruction and interval diagnostics
in this revision.
Earlier mathematical sources are retained with their original provenance.
The named author is responsible for checking all proofs and references
before submission. The accompanying finite tests and compilation records
are reproducibility evidence, not independent proof certification,
physical sensor validation or an editorial decision.
\clearpage
\begin{thebibliography}{99}\raggedright
\bibitem{BCSS}
P. Bose, J.-L. De Carufel, A. Shaikhet, and M. Smid,
\emph{Probing convex polygons with a wedge}, arXiv:1506.02572v3 (2016).
\bibitem{DGSS}
N. Dolbilin, A. Garber, E. Schulte, and M. Senechal,
\emph{Bounds for the regularity radius of Delone sets},
Discrete Comput. Geom. \textbf{74} (2025), 78--94.
\href{https://doi.org/10.1007/s00454-024-00666-6}{doi:10.1007/s00454-024-00666-6}.
\bibitem{ES}
H. Edelsbrunner and S. S. Skiena, \emph{Probing convex polygons with
X-rays}, SIAM J. Comput. \textbf{17} (1988), 870--882.
\href{https://doi.org/10.1137/0217054}{doi:10.1137/0217054}.
\bibitem{Galerne}
B. Galerne, \emph{Computation of the perimeter of measurable sets via
their covariogram. Applications to random sets}, Image Anal. Stereol.
\textbf{30} (2011), 39--51.
\href{https://doi.org/10.5566/ias.v30.p39-51}{doi:10.5566/ias.v30.p39-51}.
\bibitem{HK}
P. Herva and J. Kari, \emph{Periodicity and local complexity of Delone
sets}, arXiv:2504.20709v1 (2025).
\bibitem{LP}
J. C. Lagarias and P. A. B. Pleasants, \emph{Local complexity of Delone
sets and crystallinity}, arXiv:math/0105088v1 (2001).
\bibitem{Matheron}
G. Matheron, \emph{Random sets theory and its applications to stereology},
J. Microscopy \textbf{95} (1972), 15--23.
\href{https://doi.org/10.1111/j.1365-2818.1972.tb03708.x}{doi:10.1111/j.1365-2818.1972.tb03708.x}.
\bibitem{Molchanov}
I. Molchanov, \emph{Theory of Random Sets}, second edition,
Springer, London, 2017.
\href{https://doi.org/10.1007/978-1-4471-7349-6}{doi:10.1007/978-1-4471-7349-6}.
\bibitem{Supplement}
Q. Qi, \emph{Reference-free certification from intrinsic boundary laws},
A2 v28, author revision of October 3, 2026. Complete reviewed source
included as Supplement P, nested in the retained v29 submission;
not cited as a journal publication.
\bibitem{LawlerLimic}
G. F. Lawler and V. Limic, \emph{Random Walk: A Modern Introduction},
Cambridge Studies in Advanced Mathematics \textbf{123}, Cambridge
University Press, 2010.
\bibitem{Bertsekas}
D. P. Bertsekas, \emph{Proper policies in infinite-state stochastic
shortest path problems}, arXiv:1711.10129 (2017).
\end{thebibliography}

\end{document}
